\documentclass{book}
\usepackage{ctex}
\usepackage[cmex10]{amsmath}
\raggedbottom

% hyperref 放在末尾，生成PDF侧边书签 + 内部跳转链接
\usepackage{hyperref}
\hypersetup{
    colorlinks = true,
    linkcolor  = blue,
    citecolor  = blue,
    urlcolor   = blue,
    pdftitle   = {INTP+Libra Self Notes},
    pdfauthor  = {Skixkk},
    pdfpagemode=UseOutlines, % 打开PDF时自动展开侧边书签面板
}

\begin{document}
\tableofcontents
\clearpage

\chapter{分析性格}

\section{财务与人生准则的相关性和对人生选择及方向的影响}

\begin{itemize}
    \item 技术学的好，要饭要到老。投资做好，不动才能饱；爱自己，她会自己动。
    \item 金钱情况就和垃圾桶一样，都是隐私，对任何人都要含糊其辞。知道了你的金钱情况他就想好怎么给你花掉了，穷人见不得钱，他们只会消费，是消费主义的奴隶，也是性的奴隶；
    \item 不要加你对象那边的亲戚的微信，包括公公婆婆，丈母娘老丈人，别联系就是最好的关系；
    \item 面对亲戚就是说惨和表现穷；
    \item 面对发小，就是夸他过的好，活得通透；
    \item 陌生人就展现魅力和地位；
    \item 面对熟人就是装唐，他好跟你借钱；
    \item 情侣矛盾自己解决，今天吵架，明天又原谅他/她了；父母介入，性质就变为了家庭对抗；
    \item 见面热情的和见面冷淡的，与人品没有强相关性；
    \item 远离任何交际花，社交小王子，这种人底线都比较低，人只有足够假，才能八面玲珑；
    \item 任何社交都是没有义的，任何朋友都是酒肉朋友，你过得好他拖累你，你过得不好他还哄你；什么不是酒肉朋友，你的商业伙伴，你的上下游；
    \item 统一回答就是3千到四千，年龄大一点，统一回答就是五千，你就说没有存款，月光，房租，算一算，正好花完；
    \item 你如果有个小兼职，小外快，一定要守口如瓶，千万不要装逼，他真能你啊；
    \item 你练出一身小肌肉，你就说随便练了练，你就说吃蛋白粉吃的，你别说实话，没有人会盼你好的，都会阴阳你；
    \item 独立，吃喝，牙刷，牙杯子；
    \item 优越，松弛，挺拔
\end{itemize}

\subsection{财务观念对性格及决策的影响}

\section{吸引 INTP}

\subsection{intp 性格分析}

intp：温暖的机器人；疏离型依恋，在某些情况下会滑向恐惧型

功能栈（feature stack）：

\begin{itemize}
    \item 主导功能 Ti 内向思考 Ti introverted thinking
    \item 辅助功能 Ne 外向直觉
    \item 第三功能 Si 内向感觉
    \item 劣势功能 Fe 外向情感
\end{itemize}

INTP 的一个核心的心理困境
主导功能 Ti 天然地依赖自己的逻辑框架来解释一切
但是 劣势功能 Fe 在人际情感这个领域迟钝又脆弱，一个用逻辑去处理情感的人，注定在情感的世界里感到挫败，
久而久之，他们就会产生一种自我保护的策略，既然我不擅长，那我就退出，我就不投入，不暴雷短板，不靠近
就不会受伤，这就是疏离型依恋的典型形成路径。INTP 不是不需要链接，而是TA们在情感领域，受到的挫折太大
了，然后让他们自己说服自己，我不需要，只有理解这一层，才能看懂每一个mbti人格都是不完整的人类性格的偏向
intp 的不完整人类性格就是情感漠视，和情感需求暴雷恐惧，因为TA 们的情感需求长期被漠视，导致其情感需求被
掩埋在心底，只能在独处的时候去消化。INTP成为一个完整的人类，就是要补齐这一点，自己养自己，把自己养好。

这种性格缺陷的人，会被什么人吸引，很简单，同性相吸和异性相吸。INTP 不是在找恋人，也不是在找玩伴，而是
在找一个能让TA们安全走出自己的逻辑堡垒的人。首先，能穿透逻辑防御，有不触碰到需求暴露羞耻，因为他不是
一个完整的人。分析质疑解构去测试你，你能不能在思维上与TA交锋，看你逻辑上有没有漏洞，看你是不是一个讲
道理的人。INTP 相处舒服的状态在于，他们是在找一个能在逻辑交锋并不把逻辑交锋当成个人攻击的人

\section{全人形象}

穿搭、阳光、理解、品味、审美、智商、羞涩；

不展示余额、冲动、依赖、粘人、完整目的、缺一不可

只撩不追，早出晚归

\section{女性两性关系行为分析与实证}

一个女人真正喜欢一个人的时候，是会主动对你好，全心为你付出的，这份喜欢里带着慢慢的母性，
她不一定有多少积蓄，但一定会在生活的方方面面，大大小小的细节里关心你，惦记你，时不时就想
给你买点好吃的，逛街就想着给你添件新衣服。

现实中的中国规训男，就是一台垃圾的赚钱机器，压根没享受过这些，却还在自欺欺人。假装自己正被爱着\newpage

\section{跟健康的人谈恋爱}

跟健康的人谈恋爱，你会觉得他不爱你。心理机制长期浸泡在警觉里的人，对风平浪静有一种天然的排斥，
如果从小到大习惯了在叹气声、关门声和忽冷忽热的脸色里茶饮观色，神经系统就会把这种持续的紧绷，
与战栗当成爱与被爱的唯一凭证。那些在动荡里长出的心智，往往需要巨大的压迫感来确认自己的存在，
没有心惊肉跳的试探，没有精疲力尽的对峙，整个人就不知该把手脚往哪里放，外界只要稍微安静一点，
心里就会产生一种接近窒息的失重感。一个心智成熟的成年人给出的爱，本质上就是一种淡若水的能量，
对方有自己的工作要干，有自己的书要读，有自己的朋友要见，有自己的事业和事情要做，对方的生命
不是围绕着另一个人团团转的。穷怕了的人不敢存大钱，挨惯了打的人不敢睡安稳觉。

\section{经济独立的责任转移}

经济独立需要你的自我养育责任转向自己，这个时候购买任何物品，将变成了必需品与总资产或现金流的比率对比。

\section{行为分析}

一个天天找他人聊天的人，她必定没有要做的事情去做，也就是一个没有未来的人。

当一个男性天天找一个女性聊天，这个女性只会认为自己魅力大；一个男性每天跟不同的女性交流，那对你有意思
的女生会着急，也会权衡利弊，反正你捏住她了。这种天天找你聊天的女生会比较清闲。没有什么内涵，也没有什
么不可得性。男生不会爱上天天陪他聊天的女人，只会爱上有自己独立事业的女人。要是对方对你冷淡敷衍，消息
不回，你就快速切割远离对方，或者吊着她就行。死缠烂打，主动找话题，这就是舔狗行为，只会被对方轻视反感。
女生是靠你做自己，自身充满的魅力吸引来的，爱人先爱己，我的爱爱完自己，基本就没工夫爱别人了，这就是积
攒自身实力与未来，这就是吸引力。千万别搞反了，一味的讨好付出，自己却在原地踏步只会成为路边一条。

鹤立鸡群，鹤跟鸡维持什么良好关系啊。鹤跟鹤玩，鸡跟鸡玩。

\section{劣质行为}

听歌、看电视剧、看电影、读诗词古籍，就是在接受教育。而且这些东西教的更详细，直接告诉你，遇到什么情况
应该怎么做，这直接生产了一代的恋爱脑和无数错误的社会博弈观念。在现代人们看到大量的生活，是来自于各种
影视平台和流媒体平台刻意营造的虚拟剧情，目的是宣传消费主义与促进消费，为了让观众觉得好看，觉得爽，这
天然就脱离了真实的生存和生活实践。而人的大脑没有同步进化，依旧进行行为模仿进行行为行动，并没有太大的
意识分辨。这会群体性的导致大量的人在现实生活中的生存能力还是比较低级的时候，并与欲望不匹配造成的焦虑
引发了一系列的个体和群体的精神和社会问题。总结，每一个现代人都应该把娱乐内容看成是糖水烟酒来对待，能
不碰就不碰。这是偶尔的大额消费，不是常态化的超过现金流支出水平的长期支出。否则就会负债、焦虑、好高骛
远、不切实际和无意识消费与自身经济不匹配的消费品。必需品（也就是刚需）是七宗罪：傲慢、嫉妒、愤怒、懒
惰、贪婪、暴食、色欲。傲慢的需求，愤怒的需求，懒惰的需求，贪婪的需求... 这些需求要做到满足后才能成为
一个正常的人，也就是天然的远离，或者无意识自然的克制。

\section{底层人的劣根性与防范}

什么是底层人，默认所有人。

一部分老百姓

人的每一句话都是有目的的，这种人长期哭穷、卖惨、装可怜，就是说明这个东西对他有利，可以获得群体怜悯情感
、资源倾斜和道德正义，可以满足他的贪婪和自私。当你对他动了恻隐之心的时候，你就成为了他得到利益的俘虏。

\section{阶层词用意识对行为决策的影响}

蠢人认为善的东西，在聪明人那叫蠢
基于概率与科学分析的决策做出的大恩大义的行为，在穷人那里被说成冷漠自私

\section{传统及伪传统在现代社会的局限错误指导性}

传统习俗以前都是给奴隶用的。中国自改革开放以来及民国时期期间的部分城市的其他时期及地区，都是发展极其
缓慢甚至是倒退的，这导致了人际资源与物质资源的分配及代际传递是不足的，所以设计了分配制度——长幼尊卑。
来分配资源及金钱。

真正应该尊重的从来不是一个人的出生年份，而是他的知识经验、判断和持续学习的能力。现代社会更需要的是谁
掌握最新的信息，谁掌握了稀缺信息及实践经验，谁就有资格参与决策。当我研究的是一个你从来没有接触过的新
的领域，那么你过去几十年的经验并不能自动替代我今天掌握的新知识。尊重老者是社会美德，把年龄本身当成权
威，就会成为可能性的束缚。很多所谓的传统，并没有任何第一性原理的推理，而是过去人为了解决过去问题，设
计出来的一套办法。

\chapter{分析个体经济}

杠铃策略是一种“两端重、中间轻”的资产配置方法：大部分资金投向极低风险资产（如国债、现金），小部分投向
高风险高收益资产（如科技股、期权），由塔勒布在《反脆弱》中系统阐述。

拒绝预测，不预测未来涨跌，不赌市场方向，只利用收益分布的不对称性。

保守端只看本金是否安全，不看收益高低。进取端重仓单一标的，获取正向黑天鹅收益。保守端解决的是不出局，
而不是保证盈利，这就是由内发出的耐心、克制。

看看教女生如何追男生的视频，就知道了，大家都在胡说八道，本身是想买课和收门槛费、咨询费。MBTI 这种
也是资本宣传后的东西，之前火的星座，塔罗，从古至今的算卦，大六壬，梅花易数，紫微斗数是一类产品，可
以用这些和对方找话题增进感情用的。人不能简单的分类因为成长的很多经历不太一致，都有细微的差异和雷点。
社科理论只能解释过去的行为和经济，并不能预测未来。数据挖掘这种东西通常使用概率论和数理统计、高等概
率论这种工具去分析。这个时候就知道了，跑的两次结果是不一样的，也是呈现不同分布的，细微的初始条件也
不一样，导致结果也没有没有可比关系。这个时候就意识到了人与人是不一样的。通用的理论通常建立在人群分
类上，这是统计结果，通常有系统误差和例外。

\end{document}
